\ifdefined\gkver\else\def\gkver{student}\fi
\documentclass[\gkver]{gaokaozhenti}
\begin{document}

\chapter{2010年江苏}

\begin{enumerate}
\item
%% number: 1
%% paperName: 2010年江苏高考第 1 题 3 分
%% typeId: 1
%% type: 单选题
%% chapter: 直线运动
%% point: 运动的合成
%% method: 无
%% score: 3
%% degree: 650
%% duplicateId: 0
%% body:
如图所示，一块橡皮用细线悬挂于 O 点，用铅笔靠着线的左侧水平向右匀速移动，运动中始终保持悬线竖直，则橡皮运动的速度\xzanswer{A}

\onepicture[w=4.5cm]{figs/01.png}

\fourchoices[answer=A]
{大小和方向均不变}
{大小不变，方向改变}
{大小改变，方向不变}
{大小和方向均改变}

%% answer: A
%% memo:
\memoanswer{
橡皮在水平方向匀速运动，在竖直方向匀速运动，合运动是匀速运动。
}

\item
%% number: 2
%% paperName: 2010年江苏高考第 2 题 3 分
%% typeId: 1
%% type: 单选题
%% chapter: 电磁感应
%% point: 法拉第电磁感应定律
%% method: 无
%% score: 3
%% degree: 650
%% duplicateId: 0
%% body:
一矩形线框置于匀强磁场中，线框平面与磁场方向垂直，先保持线框的面积不变，将磁感应强度在 $1\Us$ 时间内均匀地增大到原来的两倍，接着保持增大后的磁感应强度不变，在 $1\Us$ 时间内，再将线框的面积均匀地减小到原来的一半，先后两个过程中，线框中感应电动势的比值为\xzanswer{B}

\fourchoices[answer=B]
{$\frac{1}{2}$}
{1}
{2}
{4}

%% answer: B
%% memo:
\memoanswer{
由法拉第电磁感应定律 $E=\frac{\Delta\Phi}{\Delta t}$，且 $\Delta\Phi=\Delta BS$、$\Delta\Phi=B\Delta S$，有 $E_{1}=\frac{\Delta BS}{t}=\frac{(2B-B)S}{t}=\frac{BS}{t}$，$E_{2}=\frac{2B\Delta S}{t}=\frac{2B\left(\frac{1}{2}S-S\right)}{t}=-\frac{BS}{t}$，大小相等。选项 B 正确。
}

\item
%% number: 3
%% paperName: 2010年江苏高考第 3 题 3 分
%% typeId: 1
%% type: 单选题
%% chapter: 力物体的平衡
%% point: 多力平衡
%% method: 无
%% score: 3
%% degree: 450
%% duplicateId: 0
%% body:
如图所示，置于水平地面的三脚架上固定着一质量为 $m$ 的照相机，三脚架的三根轻质支架等长，与竖直方向均成 $30^{\circ}$ 角，则每根支架中承受的压力大小为\xzanswer{D}

\onepicture[w=3.5cm]{figs/03.png}

\fourchoices[answer=D]
{$\frac{1}{3}mg$}
{$\frac{2}{3}mg$}
{$\frac{\sqrt{3}}{6}mg$}
{$\frac{2\sqrt{3}}{9}mg$}

%% answer: D
%% memo:
\memoanswer{
由力的合成及平衡可得 $3F\cos30^{\circ}=mg$，$F=\frac{2\sqrt{3}}{9}mg$，选项 D 正确。
}

\item
%% number: 4
%% paperName: 2010年江苏高考第 4 题 3 分
%% typeId: 1
%% type: 单选题
%% chapter: 交流电
%% point: LC振荡回路
%% method: 无
%% score: 3
%% degree: 650
%% duplicateId: 0
%% body:
如图所示的电路中，电源的电动势为 $E$，内阻为 $r$，电感 $L$ 的电阻不计，电阻 $R$ 的阻值大于灯泡 D 的阻值，在 $t=0$ 时刻闭合开关 S，经过一段时间后，在 $t=t_{1}$ 时刻断开 S，下列表示 A、B 两点间电压 $U_{AB}$ 随时间 $t$ 变化的图象中，正确的是\xzanswer{B}

\onepicture[w=4.5cm]{figs/04a.png}

\fourchoices[answer=B, ispicture=true, h=2.4cm]
{figs/04b.png}
{figs/04c.png}
{figs/04d.png}
{figs/04e.png}

%% answer: B
%% memo:
\memoanswer{
开关闭合时，线圈由于自感对电流的阻碍作用，可看做电阻，线圈电阻逐渐减小，并联电路电阻逐渐减小，电压 $U_{AB}$ 逐渐减小；开关闭合后再断开时，线圈的感应电流与原电流方向相同，形成回路，灯泡的电流与原电流方向相反，并逐渐减小到 0，所以正确选项 B。
}

\item
%% number: 5
%% paperName: 2010年江苏高考第 5 题 3 分
%% typeId: 1
%% type: 单选题
%% chapter: 电场
%% point: 电场中的图象
%% method: 无
%% score: 3
%% degree: 450
%% duplicateId: 0
%% body:
空间有一沿 $x$ 轴对称分布的电场，其电场强度 $E$ 随 $x$ 变化的图象如图所示。下列说法正确的是\xzanswer{C}

\onepicture[w=6.0cm]{figs/05.png}

\fourchoices[answer=C]
{O 点的电势最低}
{$x_{2}$ 点的电势最高}
{$x_{1}$ 和 $-x_{1}$ 两点的电势相等}
{$x_{1}$ 和 $x_{3}$ 两点的电势相等}

%% answer: C
%% memo:
\memoanswer{
沿 $x$ 轴对称分布的电场，由题图可得其电场线以 O 点为中心指向正、负方向，沿电场线电势降落（最快），所以 O 点电势最高，A 错误，B 错误；根据 $U=Ed$，电场强度是变量，可用 $E-x$ 图象面积表示，所以 C 正确；两点电场强度大小相等，电势不相等，D 错误。
}

\item
%% number: 6
%% paperName: 2010年江苏高考第 6 题 4 分
%% typeId: 2
%% type: 多选题
%% chapter: 曲线运动
%% point: 变轨问题
%% method: 无
%% score: 4
%% degree: 450
%% duplicateId: 0
%% body:
2009 年 5 月，航天飞机在完成对哈勃空间望远镜的维修任务后，在 A 点从圆形轨道Ⅰ进入椭圆轨道Ⅱ，B 为轨道Ⅱ上的一点，如图所示，关于航天飞机的运动，下列说法中正确的有\xzanswer{ABC}

\onepicture[w=4.5cm]{figs/06.png}

\fourchoices[answer=ABC]
{在轨道Ⅱ上经过 A 的速度小于经过 B 的速度}
{在轨道Ⅱ上经过 A 的动能小于在轨道Ⅰ上经过 A 的动能}
{在轨道Ⅱ上运动的周期小于在轨道Ⅰ上运动的周期}
{在轨道Ⅱ上经过 A 的加速度小于在轨道Ⅰ上经过 A 的加速度}

%% answer: ABC
%% memo:
\memoanswer{
根据开普勒定律，近地点的速度大于远地点的速度，A 正确。由Ⅰ轨道变到Ⅱ轨道要减速，所以 B 正确。根据开普勒定律，$\frac{R^{3}}{T^{2}}=c$，$R_{\mathrm{II}}<R_{\mathrm{I}}$，所以 $T_{\mathrm{II}}<T_{\mathrm{I}}$，C 正确。根据 $G\frac{Mm}{R^{2}}=ma$ 得 $a=\frac{GM}{R^{2}}$，又 $R_{\mathrm{II}}<R_{\mathrm{I}}$，所以 $a_{\mathrm{II}}>a_{\mathrm{I}}$，D 错误。
}

\item
%% number: 7
%% paperName: 2010年江苏高考第 7 题 4 分
%% typeId: 2
%% type: 多选题
%% chapter: 电磁感应
%% point: 远距离输电
%% method: 无
%% score: 4
%% degree: 650
%% duplicateId: 0
%% body:
在如图所示的远距离输电电路图中，升压变压器和降压变压器均为理想变压器，发电厂的输出电压和输电线的电阻均不变，随着发电厂输出功率的增大，下列说法中正确的有\xzanswer{CD}

\onepicture[w=6.5cm]{figs/07.png}

\fourchoices[answer=CD]
{升压变压器的输出电压增大}
{降压变压器的输出电压增大}
{输电线上损耗的功率增大}
{输电线上损耗的功率占总功率的比例增大}

%% answer: CD
%% memo:
\memoanswer{
变压器作用是变压，发电厂的输出电压不变，升压变压器输出的电压 $U_{2}$ 应不变，A 错误。

由于输电线电流 $I=\frac{P}{U_{2}}$，输电线电压损失 $U_{\text{损}}=IR$，降压变压器的初级电压 $U_{3}=U_{2}-U_{\text{损}}$，因 $P$ 变大，$I$ 变大，所以 $U_{\text{损}}$ 变大，降压变压器初级电压 $U_{3}$ 变小，输出电压 $U_{4}$ 也变小，B 错误。

输电线功率损失 $P_{\text{损}}=I^{2}R$，因 $I$ 变大，所以 $P_{\text{损}}$ 变大，C 正确。

$\frac{P_{\text{损}}}{P}=\frac{I^{2}R}{P}=\frac{PR}{U_{2}^{2}}$，因 $P$ 变大，所以比值变大，D 正确。
}

\item
%% number: 8
%% paperName: 2010年江苏高考第 8 题 4 分
%% typeId: 2
%% type: 多选题
%% chapter: 功和能
%% point: 动能定理
%% method: 无
%% score: 4
%% degree: 250
%% duplicateId: 0
%% body:
如图所示，平直木板 AB 倾斜放置，板上的 P 点距 A 端较近，小物块与木板间的动摩擦因数由 A 到 B 逐渐减小，先让物块从 A 由静止开始滑到 B。然后，将 A 着地，抬高 B，使木板的倾角与前一过程相同，再让物块从 B 由静止开始滑到 A。上述两过程相比较，下列说法中一定正确的有\xzanswer{AD}

\onepicture[w=4.5cm]{figs/08.png}

\fourchoices[answer=AD]
{物块经过 P 点的动能，前一过程较小}
{物块从顶端滑到 P 点的过程中因摩擦产生的热量，前一过程较少}
{物块滑到底端的速度，前一过程较大}
{物块从顶端滑到底端的时间，前一过程较长}

%% answer: AD
%% memo:
\memoanswer{
由于 $a=g\sin\alpha-\mu g\cos\alpha$，动摩擦因数 $\mu$ 由 A 到 B 逐渐减小，两过程倾角相同，前一过程加速度 $a_{1}$ 逐渐增大，后一过程加速度 $a_{2}$ 逐渐减小，即 $a_{1}<a_{2}$。根据 $v^{2}=2as$，P 点距 A 端较近，$E_{k}=\frac{1}{2}mv^{2}$，所以 $E_{k1}<E_{k2}$，选项 A 正确。

由于摩擦产生的热量 $Q=\mu mg\cos\alpha\cdot s$，且 $\mu_{1}>\mu_{2}$，P 点距 A 端较近，无法确定两过程中 $Q$ 的大小，选项 B 错误。

因 $W_{f}=\mu mg\cos\alpha\cdot s$，动摩擦因数 $\mu$ 由 A 到 B 或由 B 到 A 变化相同，$W_{f}$ 相同，由动能定理 $\frac{1}{2}mv^{2}=mgh-W_{f}$，所以物块滑到底端的速度相同，C 错误。

由于前一过程加速度 $a_{1}$ 逐渐增大，后一过程加速度 $a_{2}$ 逐渐减小，物块从顶端滑到底端的时间，前一过程较长，选项 D 正确。
}

\item
%% number: 9
%% paperName: 2010年江苏高考第 9 题 4 分
%% typeId: 2
%% type: 多选题
%% chapter: 磁场
%% point: 带电粒子在磁场中的运动
%% method: 无
%% score: 4
%% degree: 250
%% duplicateId: 0
%% body:
如图所示，在匀强磁场中附加另一匀强磁场，附加磁场位于图中阴影区域，附加磁场区域的对称轴 OO$'$ 与 SS$'$ 垂直。$a$、$b$、$c$ 三个质子先后从 S 点沿垂直于磁场的方向射入磁场，它们的速度大小相等，$b$ 的速度方向与 SS$'$ 垂直，$a$、$c$ 的速度方向与 $b$ 的速度方向间的夹角分别为 $\alpha$、$\beta$，且 $\alpha>\beta$。三个质子经过附加磁场区域后能达到同一点 S$'$，则下列说法中正确的有\xzanswer{CD}

\onepicture[w=5.5cm]{figs/09.png}

\fourchoices[answer=CD]
{三个质子从 S 运动到 S$'$ 的时间相等}
{三个质子在附加磁场以外区域运动时，运动轨迹的圆心均在 OO$'$ 轴上}
{若撤去附加磁场，$a$ 到达 SS$'$ 连线上的位置距 S 点最近}
{附加磁场方向与原磁场方向相同}

%% answer: CD
%% memo:
\memoanswer{
由于在匀强磁场中的路径大小不同，三个质子从 S 运动到 S$'$ 的时间不相等，A 错误。三个质子在附加磁场中只有 $b$ 的运动轨迹的圆心在 OO$'$ 轴上，所以 B 错误。撤去附加磁场由作图法可得，$a$ 到达 SS$'$ 连线的位置距 S 最近，选项 C 正确。三个质子在原磁场的曲率半径是相同的，无附加磁场时 $a$ 到达 SS$'$ 连线的位置距 S 最近，要使三个质子经过附加磁场区域后能达到同一点 S$'$，附加磁场方向与原磁场方向相同，选项 D 正确。
}

\item
%% number: 10
%% paperName: 2010年江苏高考第 10 题 8 分
%% typeId: 4
%% type: 实验题
%% chapter: 电路
%% point: 测电动势与内阻
%% method: 无
%% score: 8
%% degree: 450
%% duplicateId: 0
%% body:
在测量电源的电动势和内阻的实验中，由于所用的电压表（视为理想电压表）的量程较小，某同学设计了如图所示的实物电路。
\begin{enumerate}
	\item
	实验时，应先将电阻箱的电阻调到\tkanswer[2.0cm]{最大值}（选填“最大值”、“最小值”或“任意值”）。
	\item
	改变电阻箱的阻值 $R$，分别测出阻值 $R_{0}=10\UO$ 的定值电阻两端的电压 $U$，下列两组 $R$ 的取值方案中，比较合理的方案是\tkanswer[1.5cm]{2}（选填 1 或 2）。

	\begin{table}[htbp]
	\centering
	\begin{tblr}{|c|c|c|c|c|c|}
	\hline
	方案编号 & \SetCell[c=5]{c} 电阻箱的阻值 $R/\UO$ & & & & \\
	\hline
	1 & 400.0 & 350.0 & 300.0 & 250.0 & 200.0 \\
	\hline
	2 & 80.0 & 70.0 & 60.0 & 50.0 & 40.0 \\
	\hline
	\end{tblr}
	\end{table}
	\item
	根据实验数据描点，绘出的 $\frac{1}{U}-R$ 图象是一条直线。若直线的斜率为 $k$，在 $\frac{1}{U}$ 坐标轴上的截距为 $b$，则该电源的电动势 $E=$\tkanswer[2.2cm]{$\frac{1}{kR_{0}}$}，内阻 $r=$\tkanswer[2.6cm]{$\frac{b}{k}-R_{0}$}（用 $k$、$b$ 和 $R_{0}$ 表示）。
\end{enumerate}

\onepicture[w=7.0cm, align=right]{figs/10.png}

%% answer:
\jdanswer{
\begin{enumerate}
	\item
	最大值

	\item
	2

	\item
	$\frac{1}{kR_{0}}$，$\frac{b}{k}-R_{0}$
\end{enumerate}
}
%% memo:
\memoanswer{
（1）电阻箱接入回路中的电阻越大，回路电流越小，起保护电源、防止烧坏电压表的作用。

（2）方案 1 中电阻箱接入电路的电阻太大，电压表示数太小，误差太大，因此采用方案 2。或若 $R=300\UO$，电流约为 $\frac{9}{300+10}\UA=0.03\UA$；若 $R=60\UO$，电流约为 $\frac{9}{60+10}\UA=0.15\UA$，后者电流大误差小，所以填 2。

（3）电流 $I=\frac{U}{R_{0}}$，$E=I(R+R_{0}+r)$，得 $\frac{1}{U}=\frac{1}{ER_{0}}R+\frac{1}{E}+\frac{r}{ER_{0}}$。$\frac{1}{U}-R$ 图象的斜率 $k=\frac{1}{ER_{0}}$，截距 $b=\frac{1}{E}+\frac{r}{ER_{0}}$，所以电动势 $E=\frac{1}{kR_{0}}$，内阻 $r=\frac{b}{k}-R_{0}$。
}

\item
%% number: 11
%% paperName: 2010年江苏高考第 11 题 10 分
%% typeId: 4
%% type: 实验题
%% chapter: 运动和力
%% point: 验证牛顿第二定律
%% method: 无
%% score: 10
%% degree: 450
%% duplicateId: 0
%% body:
为了探究受到空气阻力时，物体运动速度随时间的变化规律，某同学采用了“加速度与物体质量、物体受力关系”的实验装置（如图所示）。实验时，平衡小车与木板之间的摩擦力后，在小车上安装一薄板，以增大空气对小车运动的阻力。
\begin{enumerate}
	\item
	往砝码盘中加入一小砝码，在释放小车\tkanswer[1.5cm]{之前}（选填“之前”或“之后”）接通打点计时器的电源，在纸带上打出一系列的点。
	\item
	从纸带上选取若干计数点进行测量，得出各计数点的时间 $t$ 与速度 $v$ 的数据如下表：

	\begin{table}[htbp]
	\centering
	\begin{tblr}{|c|c|c|c|c|c|c|}
	\hline
	时间 $t/\Us$ & 0 & 0.50 & 1.00 & 1.50 & 2.00 & 2.50 \\
	\hline
	速度 $v/(\Umdsn)$ & 0.12 & 0.19 & 0.23 & 0.26 & 0.28 & 0.29 \\
	\hline
	\end{tblr}
	\end{table}

	请根据实验数据作出小车的 $v-t$ 图象。

	\onepicture[w=6.0cm, align=right]{figs/11b.png}
	\item
	通过对实验结果的分析，该同学认为：随着运动速度的增加，小车所受的空气阻力将变大，你是否同意他的观点？请根据 $v-t$ 图象简要阐述理由。
\end{enumerate}

\onepicture[w=7.0cm, align=right]{figs/11a.png}

%% answer:
\jdanswer{
\begin{enumerate}
	\item
	之前

	\item
	图略（描点连线，用光滑的曲线作图）

	\item
	同意。在 $v-t$ 图象中，速度越大时加速度越小，小车受到的合力越小，则小车受空气阻力越大。
\end{enumerate}
}
%% memo:
\memoanswer{
（1）操作规程先接通电源再释放小车。

（2）描点用光滑的曲线作图（如图）。

\onepicture[w=6.0cm, align=right]{figs/11_an.png}

（3）同意。在 $v-t$ 图象中，速度越大时加速度越小，小车受到的合力越小，则小车受空气阻力越大。
}

\item
%% number: 12
%% paperName: 2010年江苏高考第 12 题 5 分
%% typeId: 8
%% type: 简答题
%% chapter: 光学
%% point: 光电效应
%% method: 无
%% score: 5
%% degree: 650
%% duplicateId: 0
%% body:
本题为选做题，包括 A、B、C 三小题，请选定其中两题，并在相应的答题区域内作答。若三题都做，则按 A、B 两题评分。

\textbf{A．}（选修模块 3-3）（12 分）

\begin{enumerate}
	\item
	为了将空气装入气瓶内，现将一定质量的空气等温压缩，空气可视为理想气体。下列图象能正确表示该过程中空气的压强 $p$ 和体积 $V$ 关系的是\xzanswer{B}

	\fourchoices[answer=B, ispicture=true, h=2.2cm]
	{figs/12a.png}
	{figs/12b.png}
	{figs/12c.png}
	{figs/12d.png}
	\item
	在将空气压缩装入气瓶的过程中，温度保持不变，外界做了 $24\UkJ$ 的功。现潜水员背着该气瓶缓慢地潜入海底，若在此过程中，瓶中空气的质量保持不变，且放出了 $5\UkJ$ 的热量。在上述两个过程中，空气的内能共减小\tkanswer[1.4cm]{5} $\UkJ$，空气\tkanswer[1.6cm]{放出}（选填“吸收”或“放出”）总热量为\tkanswer[1.4cm]{29} $\UkJ$。
	\item
	已知潜水员在岸上和海底吸入空气的密度分别为 $1.3\Ukgmc$ 和 $2.1\Ukgmc$，空气的摩尔质量为 $0.029\Ukgmol$，阿伏伽德罗常数 $N_{A}=6.02\times10^{23}\Umoln$。若潜水员呼吸一次吸入 $2\UL$ 空气，试估算潜水员在海底比在岸上每呼吸一次多吸入空气的分子数。（结果保留一位有效数字）
\end{enumerate}

\textbf{B．}（选修模块 3-4）（12 分）

\begin{enumerate}
	\item
	激光具有相干性好，平行度好、亮度高等特点，在科学技术和日常生活中应用广泛。下面关于激光的叙述正确的是\xzanswer{D}

	\fourchoices[answer=D]
	{激光是纵波}
	{频率相同的激光在不同介质中的波长相同}
	{两束频率不同的激光能产生干涉现象}
	{利用激光平行度好的特点可以测量月球到地球的距离}
	\item
	如图所示，在杨氏双缝干涉实验中，激光的波长为 $5.30\times10^{-7}\Um$，屏上 P 点距双缝 $S_{1}$ 和 $S_{2}$ 的路程差为 $7.95\times10^{-7}\Um$。则在这里出现的应是\tkanswer[1.8cm]{暗条纹}（选填“明条纹”或“暗条纹”）。现改用波长为 $6.30\times10^{-7}\Um$ 的激光进行上述实验，保持其他条件不变，则屏上的条纹间距将\tkanswer[1.6cm]{变宽}（选填“变宽”、“变窄”或“不变”）。

	\onepicture[w=4.2cm, align=right]{figs/12e.png}
	\item
	如图所示，一束激光从 O 点由空气射入厚度均匀的介质，经下表面反射后，从上面的 A 点射出。已知入射角为 $i$，A 与 O 相距 $l$，介质的折射率为 $n$，试求介质的厚度 $d$。

	\onepicture[w=4.2cm, align=right]{figs/12f.png}
\end{enumerate}

\textbf{C．}（选修模块 3-5）（12 分）

\begin{enumerate}
	\item
	研究光电效应电路如图所示，用频率相同、强度不同的光分别照射密封真空管的钠极板（阴极 K），钠极板发射出的光电子被阳极 A 吸收，在电路中形成光电流。下列光电流 $I$ 与 AK 之间的电压 $U_{AK}$ 的关系图象中，正确的是\xzanswer{C}

	\onepicture[w=3.6cm, align=right]{figs/12g.png}

	\fourchoices[answer=C, ispicture=true, h=2.2cm]
	{figs/12h.png}
	{figs/12i.png}
	{figs/12j.png}
	{figs/12k.png}
	\item
	钠金属中的电子吸收光子的能量，从金属表面逸出，这就是光电子。光电子从金属表面逸出的过程中，其动量的大小\tkanswer[1.5cm]{减小}（选填“增大”、“减小”或“不变”），原因是\tkanswer[4.5cm]{光电子受到金属表面层中力的阻碍作用（或需要克服逸出功）}。
	\item
	已知氢原子处在第一、第二激发态的能级分别为 $-3.4\UeV$ 和 $-1.51\UeV$，金属钠的截止频率为 $5.53\times10^{14}\UHz$，普朗克常量 $h=6.63\times10^{-34}\UJcdots$。请通过计算判断，氢原子从第二激发态跃迁到第一激发态过程中发出的光照射金属钠板，能否发生光电效应。
\end{enumerate}

%% answer:
\jdanswer{
\begin{enumerate}
	\item
	A：（1）B；（2）5，放出，29；（3）$3\times10^{22}$

	\item
	B：（1）D；（2）暗条纹，变宽；（3）$d=\frac{\sqrt{n^{2}-\sin^{2}i}}{2\sin i}l$

	\item
	C：（1）C；（2）减小，光电子受到金属表面层中力的阻碍作用（或需要克服逸出功）；（3）$E=1.89\UeV$，$W_{0}=2.3\UeV$，因为 $E<W_{0}$，所以不能发生光电效应。
\end{enumerate}
}
%% memo:
\memoanswer{
\textbf{A．}（1）由玻意耳定律 $pV=C$，即 $p\propto\frac{1}{V}$，选项 B 正确。

（2）根据热力学第一定律 $\Delta U=W+Q$，第一阶段 $W_{1}=24\UkJ$，$\Delta U_{1}=0$，所以 $Q_{1}=-24\UkJ$，放热；第二阶段 $W_{2}=0$，放热 $Q_{2}=-5\UkJ$，所以 $\Delta U_{2}=-5\UkJ$。又 $\Delta U=\Delta U_{1}+\Delta U_{2}=-5\UkJ$，即内能减少 $5\UkJ$；$Q=Q_{1}+Q_{2}=-29\UkJ$，即放出热量 $29\UkJ$。

（3）设空气的摩尔质量为 $M$，在海底和岸上的密度分别为 $\rho_{\text{海}}$ 和 $\rho_{\text{岸}}$，一次吸入空气的体积为 $V$，则有 $\Delta n=\frac{(\rho_{\text{海}}-\rho_{\text{岸}})V}{M}N_{A}$，代入数据得 $\Delta n=3\times10^{22}$。

\textbf{B．}（1）测距利用的是平行度好的特点，答案 D。

（2）由于 $\frac{7.95}{5.30}=1.5$，路程差是半波长的奇数倍，所以是暗条纹。又 $\Delta x=\frac{l}{d}\lambda$，$\lambda$ 变大，$\Delta x$ 变大，条纹变宽。

（3）设折射角为 $\gamma$，由折射定律 $\frac{\sin i}{\sin\gamma}=n$；由几何关系 $l=2d\tan\gamma$，解得 $d=\frac{\sqrt{n^{2}-\sin^{2}i}}{2\sin i}l$。

\textbf{C．}（1）虽然 $I_{\text{强}}>I_{\text{弱}}$，但截止电压相等，选项 C 正确。

（2）减小；光电子受到金属表面层中力的阻碍作用（或需要克服逸出功）。

（3）氢原子放出的光子能量 $E=E_{3}-E_{2}$，代入数据得 $E=1.89\UeV$；金属钠的逸出功 $W_{0}=h\nu_{0}$，代入数据得 $W_{0}=2.3\UeV$。因为 $E<W_{0}$，所以不能发生光电效应。
}

\item
%% number: 13
%% paperName: 2010年江苏高考第 13 题 15 分
%% typeId: 6
%% type: 计算题
%% chapter: 电磁感应
%% point: 电磁感应中的变加速
%% method: 无
%% score: 15
%% degree: 450
%% duplicateId: 0
%% body:
如图所示，两足够长的光滑金属导轨竖直放置，相距为 $L$，一理想电流表与两导轨相连，匀强磁场与导轨平面垂直。一质量为 $m$、有效电阻为 $R$ 的导体棒在距磁场上边界 $h$ 处静止释放。导体棒进入磁场后，流经电流表的电流逐渐减小，最终稳定为 $I$。整个运动过程中，导体棒与导轨接触良好，且始终保持水平，不计导轨的电阻。求：
\begin{enumerate}
	\item
	磁感应强度的大小 $B$；
	\item
	电流稳定后，导体棒运动速度的大小 $v$；
	\item
	流经电流表电流的最大值 $I_{m}$。
\end{enumerate}

\onepicture[align=right]{figs/13.png}

%% answer:
\jdanswer{
\begin{enumerate}
	\item
	$\frac{mg}{Il}$

	\item
	$\frac{I^{2}R}{mg}$

	\item
	$\frac{mg\sqrt{2gh}}{IR}$
\end{enumerate}
}
%% memo:
\memoanswer{
（1）电流稳定后，导体棒做匀速运动 $BIl=mg$ ①。

解得 $B=\frac{mg}{Il}$ ②。

（2）感应电动势 $E=Blv$ ③。

感应电流 $I=\frac{E}{R}$ ④。

由②③④解得 $v=\frac{I^{2}R}{mg}$。

（3）由题意知，导体棒刚进入磁场时的速度最大，设为 $v_{m}$。

由机械能守恒 $\frac{1}{2}mv_{m}^{2}=mgh$，得 $v_{m}=\sqrt{2gh}$。

感应电动势的最大值 $E_{m}=Blv_{m}$。

感应电流的最大值 $I_{m}=\frac{E_{m}}{R}$。

解得 $I_{m}=\frac{mg\sqrt{2gh}}{IR}$。
}

\item
%% number: 14
%% paperName: 2010年江苏高考第 14 题 16 分
%% typeId: 6
%% type: 计算题
%% chapter: 曲线运动
%% point: 平抛运动
%% method: 极值
%% score: 16
%% degree: 450
%% duplicateId: 0
%% body:
在游乐节目中，选手需要借助悬挂在高处的绳飞越到水面的浮台上，小明和小阳观看后对此进行了讨论。如图所示，他们将选手简化为质量 $m=60\Ukg$ 的质点，选手抓住绳由静止开始摆动，此时绳与竖直方向夹角 $\alpha=53^{\circ}$，绳的悬挂点 O 距水面的高度为 $H=3\Um$。不考虑空气阻力和绳的质量，浮台露出水面的高度不计，水足够深。取重力加速度 $g=10\Umsq$，$\sin53^{\circ}=0.8$，$\cos53^{\circ}=0.6$。
\begin{enumerate}
	\item
	求选手摆到最低点时对绳拉力的大小 $F$；
	\item
	若绳长 $l=2\Um$，选手摆到最高点时松手落入手中。设水对选手的平均浮力 $f_{1}=800\UN$，平均阻力 $f_{2}=700\UN$，求选手落入水中的深度 $d$；
	\item
	若选手摆到最低点时松手，小明认为绳越长，在浮台上的落点距岸边越远；小阳认为绳越短，落点距岸边越远，请通过推算说明你的观点。
\end{enumerate}

\onepicture[align=right]{figs/14.png}

%% answer:
\jdanswer{
\begin{enumerate}
	\item
	$1080\UN$

	\item
	$1.2\Um$

	\item
	两人的看法均不正确。当绳长越接近 $1.5\Um$ 时，落点距岸边越远。
\end{enumerate}
}
%% memo:
\memoanswer{
（1）机械能守恒 $mgl(1-\cos\alpha)=\frac{1}{2}mv^{2}$ ①。

圆周运动 $F'-mg=m\frac{v^{2}}{l}$ ②。

解得 $F'=mg(3-2\cos\alpha)$。

人对绳的拉力 $F=F'$。

则 $F=1080\UN$。

（2）由动能定理 $mg(H-l\cos\alpha+d)-(f_{1}+f_{2})d=0$。

则 $d=\frac{mg(H-l\cos\alpha)}{f_{1}+f_{2}-mg}$。

解得 $d=1.2\Um$。

（3）选手从最低点开始做平抛运动 $x=vt$。

$H-l=\frac{1}{2}gt^{2}$。

且由①式解得 $x=2\sqrt{l(H-l)(1-\cos\alpha)}$。

当 $l=\frac{H}{2}$ 时，$x$ 有最大值，解得 $l=1.5\Um$。

因此，两人的看法均不正确。当绳长越接近 $1.5\Um$ 时，落点距岸边越远。
}

\item
%% number: 15
%% paperName: 2010年江苏高考第 15 题 16 分
%% typeId: 6
%% type: 计算题
%% chapter: 电场
%% point: 带电粒子的偏转
%% method: 无
%% score: 16
%% degree: 250
%% duplicateId: 0
%% body:
制备纳米薄膜装置的工作电极可简化为真空中间距为 $d$ 的两平行极板，如图~\ref{2010江苏15a} 所示，加在极板 A、B 间的电压 $U_{AB}$ 作周期性变化，其正向电压为 $U_{0}$，反向电压为 $-kU_{0}$（$k>1$），电压变化的周期为 $2\tau$，如图~\ref{2010江苏15b} 所示。在 $t=0$ 时，极板 B 附近的一个电子，质量为 $m$、电荷量为 $e$，受电场作用由静止开始运动。若整个运动过程中，电子未碰到极板 A，且不考虑重力作用。
\begin{enumerate}
	\item
	若 $k=\frac{5}{4}$，电子在 $0\sim2\tau$ 时间内不能到达极板 A，求 $d$ 应满足的条件；
	\item
	若电子在 $0\sim2\tau$ 时间未碰到极板 B，求此运动过程中电子速度 $v$ 随时间 $t$ 变化的关系；
	\item
	若电子在第 $N$ 个周期内的位移为零，求 $k$ 的值。
\end{enumerate}

\twopicture[numstyle=chinese, labelA=2010江苏15a, labelB=2010江苏15b, wA=5.0cm, wB=5.0cm, align=right]
{figs/15a.png}
{figs/15b.png}

%% answer:
\jdanswer{
\begin{enumerate}
	\item
	$d>\sqrt{\frac{9eU_{0}\tau^{2}}{10m}}$

	\item
	①当 $0\leqslant t-2n\tau<\tau$ 时，$v=[t-(k+1)n\tau]\frac{ekU_{0}}{md}$［$n=0,1,2,\cdots,(N-1)$］；②当 $0\leqslant t-(2n+1)\tau<\tau$ 时，$v=[(n+1)(k+1)\tau-kt]\frac{eU_{0}}{md}$［$n=0,1,2,\cdots,(N-1)$］

	\item
	$k=\frac{4N-1}{4N-3}$
\end{enumerate}
}
%% memo:
\memoanswer{
（1）电子在 $0\sim\tau$ 时间内做匀加速运动，加速度的大小 $a_{1}=\frac{eU_{0}}{md}$，位移 $x_{1}=\frac{1}{2}a_{1}\tau^{2}$。

在 $\tau\sim2\tau$ 时间内先做匀减速运动，后反向做匀加速运动，加速度的大小 $a_{2}=\frac{5eU_{0}}{4md}$，初速度的大小 $v_{1}=a_{1}\tau$，匀减速运动阶段的位移 $x_{2}=\frac{v_{1}^{2}}{2a_{2}}$。由 $d>x_{1}+x_{2}$，解得 $d>\sqrt{\frac{9eU_{0}\tau^{2}}{10m}}$。

（2）在 $2n\tau\sim(2n+1)\tau$［$n=0,1,2,\cdots,(N-1)$］时间内速度增量 $\Delta v_{1}=a_{1}\tau$。

在 $(2n+1)\tau\sim2(n+1)\tau$［$n=0,1,2,\cdots,(N-1)$］时间内，加速度的大小 $a_{2}'=\frac{ekU_{0}}{md}$，速度增量 $\Delta v_{2}=-a_{2}'\tau$。

①当 $0\leqslant t-2n\tau<\tau$ 时，电子的运动速度 $v=n\Delta v_{1}+n\Delta v_{2}+a_{1}(t-2n\tau)$，解得 $v=[t-(k+1)n\tau]\frac{ekU_{0}}{md}$［$n=0,1,2,\cdots,(N-1)$］。

②当 $0\leqslant t-(2n+1)\tau<\tau$ 时，电子的运动速度 $v=(n+1)\Delta v_{1}+n\Delta v_{2}-a_{2}'[t-(2n+1)\tau]$，解得 $v=[(n+1)(k+1)\tau-kt]\frac{eU_{0}}{md}$［$n=0,1,2,\cdots,(N-1)$］。

（3）电子在 $2(N-1)\tau\sim(2N-1)\tau$ 时间内的位移 $x_{2N-1}=v_{2N-2}\tau+\frac{1}{2}a_{1}\tau^{2}$，电子在 $(2N-1)\tau\sim2N\tau$ 时间内的位移 $x_{2N}=v_{2N-1}\tau-\frac{1}{2}a_{2}'\tau^{2}$。

由以上两式可知 $v_{2N-2}=(N-1)(1-k)\tau\frac{eU_{0}}{md}$，$v_{2N-1}=(N-Nk+k)\tau\frac{eU_{0}}{md}$。

依题意得 $x_{2N-1}+x_{2N}=0$，解得 $k=\frac{4N-1}{4N-3}$。
}

\end{enumerate}

\end{document}
